\section{Small Fourier signals}
\label{sec:signals}

The purpose of this section is to modulate a packet's spectral density
according to its angular index.  The size of the modulation will depend on
a fixed band of \emph{normalized} indices, but not on the scale of those
indices.  The integer added to the old speed is allowed to depend on the
actual, finite set of indices.

Fix numbers $0<b<D_0$ and a function
\begin{equation}
 d\in C_c^\infty(\R\setminus\{0\};\C),
 \qquad d(-x)=\overline{d(x)}.
 \label{eq:signal-d}
\end{equation}
A scale is a positive real number $B$, and the normalized index of a
packet at angular index $m$ is $x=m/B$.

\begin{lemma}[Transmission of an index signal]
\label{lem:signal}
For the fixed data $b,D_0,d$ there are an integer $P\ge2$ and a number
$\gamma_0>0$ with the following property.  Any fixed
$0<\gamma\le\gamma_0$ can be used, independently of $B$.

Let $S$ be an intermediate map, and let
\[
 F_r=\sum_{m\in\mathcal M_r} A_{r,m}(u,a)\ee(m\theta)
\]
be finite smooth packet data in its good arc.  Suppose that the finite
sets $\mathcal M_r\subset\Z\setminus\{0\}$ are mutually disjoint,
$|A_{r,m}(u,a)|$ is independent of $a$, and
$b\le |m|/B\le D_0$.  Write $p_{r,m}$ for the spectral density of the
indicated packet. Given any $\eps,\delta_x,\delta_z>0$ and prescribed
positive bounds for the physical $L^2$ transport errors, a passage produces
finite output packet sums $F_{r,m}^+$, approximating the corresponding
input packets in physical $L^2$, whose spectral densities satisfy
\begin{equation}
 \sum_{r,m}\left\|p_{r,m}^+
 -p_{r,m}(z)\left[1+2\gamma\1_{\gcd(m,P)=1}
       \Rea\{d(m/B)\ee(-Nz)\}\right]\right\|_{L^1(\T)}
 <\eps.
 \label{eq:signal-conclusion}
\end{equation}
The new scale is $B^+=NcB$, and the new speed is $\sigma_N$ from
Lemma~\ref{lem:passage}. For every point $(u',a')$ in the coefficient
support of a retained output packet of index $j$ descending from
$(r,m)$, there is $(u,a)\in\supp A_{r,m}$ such that
\begin{equation}
 \left|\frac{j}{B^+}-\frac mB\right|<\delta_x,
 \qquad
 \dist_\T\bigl(j\sigma_N(u'),m\phi(u)\bigr)<\delta_z.
 \label{eq:signal-support-closeness}
\end{equation}
Its modulus is independent of the new spare coordinate. Distinct old
angular indices have disjoint output angular-index supports.
Conjugate symmetry is preserved when present.

The passage and its inverse can lie in any prescribed smooth
neighborhood of $S$ and $S^{-1}$ in physical coordinates, and the new
good arc can have arbitrarily small complement.  The choices can be made
in the order $c$, a bounded smooth jitter and cutoffs, and finally the
integer $N$.  After the preceding choices, every sufficiently large
integer $N$ works.  In particular an arbitrary lower bound on $N$ can
be imposed.
\end{lemma}

\begin{proof}
\emph{A positive law.}
We first construct a positive law without using $B$.  We use the
positive-exponent Fourier convention
\[
 \widehat q(v)=\int_\R q(R)\ee(vR)\,\dd R.
\]
Choose $0<\alpha<b/2$, put $a_0=(2\alpha)^{-1}$, and define
\begin{equation}
 s_\alpha(R)=\alpha\left(\frac{\sin(\pi\alpha R)}
                                      {\pi\alpha R}\right)^2,
 \qquad
 q_0(R)=\frac1{2a_0}\int_{-a_0}^{a_0}s_\alpha(R-t)\,\dd t,
 \label{eq:signal-baseline}
\end{equation}
with the removable value at the origin understood.  Fourier inversion
for an interval, followed by convolution, gives
\[
 \widehat{s_\alpha}(v)=(1-|v|/\alpha)_+.
\]
Consequently $q_0$ is a smooth probability density with Fourier support
in $[-\alpha,\alpha]$.  It is strictly positive: the interval of
translations in \eqref{eq:signal-baseline} cannot be contained in the
discrete zero set of $s_\alpha$.

Here a quantitative lower tail is useful.  Since the averaging interval
has length $1/\alpha$, it contains exactly one period of
$\sin^2(\pi\alpha(R-t))$ as a function of $t$.  For
$|R|\ge2a_0+1$ this gives
\[
 \frac1{2\pi^2\alpha(|R|+a_0)^2}
 \le q_0(R)
 \le\frac1{\pi^2\alpha(|R|-a_0)^2}.
\]
Positivity and continuity on the remaining compact interval therefore
give constants $c_0,C_0>0$ such that
\begin{equation}
 \frac{c_0}{1+R^2}\le q_0(R)\le\frac{C_0}{1+R^2}
 \qquad(R\in\R).
 \label{eq:signal-quadratic-tails}
\end{equation}

Choose $D>b$ such that $\supp d\subset[-D,D]$, and then an integer
$L>\max(2,D/b)$.  Let $P$ be the product of all primes at most $L$.
The strict inequality in this choice will allow a small dilation of
the Fourier arguments.  Define
\[
 g(R)=\int_\R d(v)\ee(-vR)\,\dd v,
 \qquad
 C_P(h)=\sum_{k\in(\Z/P\Z)^\times}\ee(-kh/P)
 \quad(0\le h<P).
\]
The function $g$ is Schwartz by repeated integration by parts.  The
symmetry in \eqref{eq:signal-d} makes it real, and
$\int g=d(0)=0$.  The number $C_P(h)$ is real because the set of units
is invariant under negation.  Set
\begin{equation}
 Q_h(R)=\frac1P\bigl[q_0(R)+\gamma C_P(h)g(R)\bigr].
 \label{eq:signal-joint-law}
\end{equation}
For $G=\sup_R(1+R^2)|g(R)|$ and $C_* =\max_h|C_P(h)|$, choose
$\gamma_0>0$ so small that
\[
 \gamma_0 C_*G\le c_0/2,
 \qquad 2\gamma_0\|d\|_\infty\le1/2;
\]
a zero factor imposes no restriction.  Then, for every permitted
$\gamma$,
\begin{equation}
 \frac{q_0(R)}{2P}\le Q_h(R)\le\frac{3q_0(R)}{2P},
 \qquad \int_\R Q_h(R)\,\dd R=\frac1P.
 \label{eq:signal-positive-law}
\end{equation}
Thus $Q_h(R)\,\dd R$ defines a probability law on
$\R\times\{0,\ldots,P-1\}$.  Finite Fourier inversion gives, for every
$v\in\R$ and $k\in\Z$,
\begin{equation}
 \sum_{h=0}^{P-1}\int_\R Q_h(R)\ee(vR+kh/P)\,\dd R
 =\widehat q_0(v)\1_{k\equiv0\ (P)}
       +\gamma d(v)\1_{\gcd(k,P)=1}.
 \label{eq:signal-joint-fourier}
\end{equation}
Indeed the sum over $h$ pairs $C_P(h)$ with the character at $k$,
leaving precisely the indicator that $k$ is a unit.  All choices up to
this point are independent of $B$, the map, and its finite packet list.

\emph{Moments and the finite sieve.}
For fixed $B$ and an integer $c>0$, consider the random variable
\begin{equation}
 J_*=(R/B+h/P)/c
 \label{eq:signal-ideal-jitter}
\end{equation}
under this law.  Its density is
\[
 \rho_{B,c}(y)=cB\sum_{h=0}^{P-1}Q_h\bigl(B(cy-h/P)\bigr).
\]
It is a strictly positive smooth probability density.  The bounds
\eqref{eq:signal-positive-law} imply quadratic upper and lower tails,
with constants now allowed to depend on $B,c$.

We calculate the periodized density of $sJ_*$.  Write $x=m/B$ and
$s=mct$.  For its positive integer Fourier moments, the exact formula is
\begin{equation}
 M_\ell(t)=\frac1P\sum_{h=0}^{P-1}\ee(\ell mt h/P)
 \bigl[\widehat q_0(\ell xt)
             +\gamma C_P(h)d(\ell xt)\bigr],\qquad \ell\ge1.
 \label{eq:signal-exact-moment}
\end{equation}
At $t=1$, \eqref{eq:signal-joint-fourier} applies.  Its baseline term
vanishes since $|\ell x|\ge b>\alpha$.  If $\ell\ge2$ and
$d(\ell x)\ne0$, then $\ell\le D/b<L$.  Some prime divisor of $\ell$
therefore divides $P$, so $\gcd(\ell m,P)>1$.  It follows that
\begin{equation}
 M_1(1)=\gamma d(x)\1_{\gcd(m,P)=1},
 \qquad M_\ell(1)=0\quad(\ell\ge2).
 \label{eq:signal-ideal-moments}
\end{equation}

For finite $c$ the discrete phases in
\eqref{eq:signal-exact-moment} need not be residue characters.  We
estimate their errors before using the sieve.  Choose
$\tau\in(1/2,1)$ so that $D/(\tau b)<L$, and take $c$ large enough
that
\[
 t_N(u)=1+\frac{\phi(u)+\Delta_N}{c}\ge\tau,
 \qquad |\Delta_N|\le1/N,
\]
on the old supports, uniformly for $N\ge1$.  The baseline term in
\eqref{eq:signal-exact-moment} then vanishes for \emph{every}
$\ell\ge1$, since $|\ell xt_N|\ge\tau b>\alpha$.
The perturbation also vanishes identically when
$\ell>D/(\tau b)$.  Hence nonintegral discrete phases cannot introduce
an unbounded collection of harmonics.  For the fixed remaining finite
set, put $C_1=P^{-1}\sum_h|C_P(h)|$.  The elementary bound
$|\ee(v)-\ee(w)|\le2\pi|v-w|$ gives
\begin{equation}
 |M_\ell(t)-M_\ell(1)|
 \le\gamma C_1\ell
 \bigl[D_0\|d'\|_\infty+2\pi|m|\|d\|_\infty\bigr]|t-1|.
 \label{eq:signal-finite-c}
\end{equation}
We may thus make the sum of all moment errors arbitrarily small by
choosing $c$ after the actual finite set of old indices.  In particular
$c$ is allowed to grow with $\max|m|$; $\gamma$ does not need to decrease.
The same choice can make the ranges $m(c+\phi(\T))$ for distinct old
indices disjoint and bounded away from zero.

To pass from moments to densities, let
\[
\mathcal P_s\rho(y)
   =\sum_{k\in\Z}\frac1{|s|}\rho((y+k)/s),\qquad y\in\T.
\]
Once $c$ has been chosen, all the relevant $s$ ranges are fixed compact
sets away from zero.  No bound on these sets uniform over future
passages is needed.
For $s$ in a fixed compact subset of $\R\setminus\{0\}$, the quadratic
upper tail of $\rho_{B,c}$ makes this a jointly continuous function of
$(s,y)$, with uniformly summable tails.  Its Fourier coefficients in
the positive-exponent convention are the moments just computed.
The compact bandwidth shows that its Fourier series is a finite sum
on the relevant branches.  Fourier uniqueness and
\eqref{eq:signal-ideal-moments}--\eqref{eq:signal-finite-c} therefore
give, uniformly on those branches and in $y$,
\begin{equation}
 \mathcal P_s\rho_{B,c}(y)
 =1+2\gamma\1_{\gcd(m,P)=1}\Rea\{d(m/B)\ee(-y)\}
       +o_{c\to\infty}(1).
 \label{eq:signal-periodic-density}
\end{equation}
The error here is for the fixed old list and scale $B$.
The negative exponent follows directly from our convention:
$\int\mathcal P_s\rho(y)\ee(\ell y)\,\dd y=M_\ell$.

\emph{Realization and support.}
We now realize this law by a passage.  Lemma~\ref{lem:law-approx}
provides a bounded smooth circle function $J$, strictly increasing on
a good arc, and a smooth cutoff $\psi$ compactly supported in that arc,
such that the subprobability law
\[
 \widetilde\rho=J_*(\psi(a)^2\,\dd a)
\]
is arbitrarily close to $\rho_{B,c}$ in total variation.  Its missing
mass is included in this error.  The good arc has as small a complement
as requested.  Since $|A_{r,m}|$ is independent of $a$, its stationary
weight is a function of $s$ alone:
\[
 w_{r,m}(s)=\frac{|A_{r,m}(u,a)|^2}{|m\phi'(u)|},
 \qquad m(c+\phi(u))=s,
\]
extended by zero.  For such a nonnegative compactly supported weight,
Lemma~\ref{lem:fast-average} and contraction of total variation under
pushforward and periodization give
\begin{align}
 &\lim_{N\to\infty}\int_\R w_{r,m}(s)
   |\mathcal P_s\widetilde\rho(Ns)
               -\mathcal P_s\rho_{B,c}(Ns)|\,\dd s \notag\\
 &\qquad=\int_\R w_{r,m}(s)\int_\T
   |\mathcal P_s\widetilde\rho(y)
               -\mathcal P_s\rho_{B,c}(y)|\,\dd y\,\dd s \notag\\
 &\qquad\le\|w_{r,m}\|_1\|\widetilde\rho-\rho_{B,c}\|_{L^1(\R)}.
 \label{eq:signal-law-comparison}
\end{align}
Here densities of subprobabilities are used in the $L^1$ norm.
The realized law has a continuous compactly supported density; the
ideal law has the uniform tail bound established above.  These facts
supply the continuity and domination required for fast averaging.
The independence of the weight from $a$ is essential in this step.

Lemma~\ref{lem:passage} identifies the output density, up to an
arbitrarily small $L^1$ error, as the folding to $\T$ of
$w_{r,m}(s)\mathcal P_s\widetilde\rho(Ns)$.
In using that lemma the rounding correction is removed at the
coefficient-mass level: its replacement of $w_{r,m}(s)$ by
$w_{r,m}(s-m\Delta_N)$ costs $o(1)$ in the summed masses.  It does
not translate the fast argument $Ns$.  Alternatively,
\eqref{eq:signal-finite-c} also holds uniformly with $t=t_N$.
Since $cm$ and $N$ are integers,
\[
 p_{r,m}(z)=\sum_{k\in\Z}w_{r,m}(z+k),
 \qquad \ee(-N(z+k))=\ee(-Nz).
\]
Now \eqref{eq:signal-periodic-density},
\eqref{eq:signal-law-comparison}, and the passage density formula
prove \eqref{eq:signal-conclusion}.  Choose $c$ first to make its
error sufficiently small, then the law approximation, and then $N$.
The finite packet list allows all errors to be summed.

For completeness we give the simultaneous support estimates, which
will also apply to neutral passages.  Put
$M=\max_{r,m}|m|$, with $M=0$ for an empty packet list.
By the support cutoff in
Lemma~\ref{lem:passage}, at each retained point of an output packet at
index $j$ there is a point $u$ in its old support for which
\begin{equation}
 j/N=m(c+\phi(u)+\Delta_N)+\eta,
 \qquad |\eta|\le\eta_0,
 \label{eq:signal-retained-support}
\end{equation}
where the image enlargement $\eta_0$ can be as small as desired.
If $x'=j/(NcB)$ and $z'=j/N+(j/N)J(a)/N$ is the new temporal
frequency, then
\begin{align}
 |x'-m/B|
 &\le\frac{D_0(\|\phi\|_\infty+1)}c+\frac{\eta_0}{cB},
 \label{eq:signal-index-drift}\\
 \dist_\T(z',m\phi(u))
 &\le\frac MN+\eta_0
  +\frac{[M(c+\|\phi\|_\infty+1)+\eta_0]\|J\|_\infty}{N}.
 \label{eq:signal-temporal-drift}
\end{align}
Choose $c$ for the first bound and index separation, then $\eta_0$
and the bounded $J$, and finally $N$ for the second bound.  A large
bound on $J$ only increases the eventual requirement on $N$.
These choices also give physical smooth closeness of the passage and
its inverse and the required $L^2$ approximations, by
Lemma~\ref{lem:passage}.

Every new coefficient descending from old index $m$ has the form
$\ee(ma_{\rm new})$ times a scalar function of the new speed
coordinate.  Because different old indices have disjoint retained
new index ranges, their spare-coordinate phases are never added at
one new index.  Thus the resulting packet moduli are independent of
the new spare coordinate.  All cutoffs can be chosen in conjugate
pairs, so conjugation is retained as well.
\end{proof}

We next remove the coprimality restriction by placing a neutral passage
before transmission.  The output residue classes of the neutral
passage, rather than the input residue classes, are the ones tested in
the transmission.

\begin{lemma}[Refresh and transmission]
\label{lem:refresh-transmit}
Fix $b,D_0,d,P,\gamma$ as in Lemma~\ref{lem:signal}, and put
\begin{equation}
 \lambda=\gamma\frac{|(\Z/P\Z)^\times|}{P}>0.
 \label{eq:signal-lambda}
\end{equation}
Suppose finite colored packet data have mutually disjoint angular-index
sets and moduli independent of the spare coordinate.  Suppose their
normalized indices lie in a fixed compact subset of
$\{x:b<|x|<D_0\}$.  Let $p_r$ be the color densities, and let
$d_r:\T\to\C$ be continuous functions such that on the support of
every input packet of color $r$,
\begin{equation}
 |d(m/B)-d_r(m\phi(u)\bmod1)|\le\omega.
 \label{eq:signal-color-match}
\end{equation}
Then a neutral passage followed by a transmission gives output colors
with
\begin{equation}
 \sum_r\left\|p_r^+
  -p_r(z)\bigl[1+2\lambda\Rea\{d_r(z)\ee(-Nz)\}\bigr]
                 \right\|_1
 \le 2\gamma\omega\sum_r\|p_r\|_1+\eps,
 \label{eq:signal-refreshed}
\end{equation}
where $\eps>0$ is arbitrary and $N$ is the second passage integer.
The physical map and inverse, the physical packet vectors, and the
normalized-index and temporal-frequency displacements can all be
controlled with arbitrary prescribed positive error bounds.  Output
packet moduli remain independent of the spare coordinate, and colors
remain disjoint in angular indices.  The final $N$ can exceed any
prescribed bound.  Conjugate-symmetric data can be treated
conjugate-symmetrically.

If different colors start in separated compact normalized-index
sets, any fixed sufficiently small enlargements of those sets can
contain their outputs.  The same conclusion holds through any fixed
finite number of such pairs, with positive margins around the band
$b<|x|<D_0$.  If the input supports satisfy
$\dist_\T(m/B,m\phi(u))\le\kappa$, the output supports satisfy the
same inequality with $\kappa$ increased by an arbitrarily small amount.
\end{lemma}

\begin{proof}
Use Lemma~\ref{lem:neutral} with residue modulus $P$.  Denote the
refreshed scale by $\widetilde B$, its packet densities by
$\widetilde p_{r,m}$, and its color-residue densities by
\[
 \widetilde p_{r,h}=\sum_{m\equiv h\ (P)}\widetilde p_{r,m}.
\]
That lemma makes
\begin{equation}
 E_r=\sum_{h=0}^{P-1}
       \|\widetilde p_{r,h}-p_r/P\|_1
 \label{eq:signal-residue-error}
\end{equation}
arbitrarily small.  It also retains disjoint output angular-index
sets and gives moduli independent of the new spare coordinate.
To control its support, first choose its $c$ as large as necessary
for \eqref{eq:signal-index-drift}; next choose its cutoff and broad
linear jitter for the residue equalization; finally choose its passage
integer for \eqref{eq:signal-temporal-drift}.  Thus equalizing residues
does not prevent either displacement from being arbitrarily small.

For a continuous function $f$ on its relevant compact domain, denote
its modulus of continuity by
$\nu_f(t)=\sup_{\dist(v,w)\le t}|f(v)-f(w)|$.
If the refresh changes normalized indices by at most $\delta_x$ and
temporal frequencies by at most $\delta_z$, its retained color
supports satisfy
\begin{equation}
 |d(m/\widetilde B)-d_r(z)|
 \le\rho_r:=\omega+\nu_d(\delta_x)+\nu_{d_r}(\delta_z).
 \label{eq:signal-refreshed-match}
\end{equation}
This follows pointwise by linking an output support point to an old
one and using \eqref{eq:signal-color-match}.  In particular the extra
terms tend to zero.  The compact margin in the normalized band lets
us apply Lemma~\ref{lem:signal} with the same $P,\gamma$ after this
refresh, although its actual indices and scale may be very large.

By that lemma the transmission density, with an arbitrarily small
summed error $T_r$ for color $r$, is
\[
 \widetilde p_r(z)+2\gamma\Rea\left\{\ee(-Nz)
    \sum_{\gcd(m,P)=1}d(m/\widetilde B)
                                   \widetilde p_{r,m}(z)\right\}.
\]
Replacing the coefficient in this sum by $d_r(z)$ costs at most
$2\gamma\rho_r\|\widetilde p_r\|_1$ in $L^1$.  Also,
\[
 \|\widetilde p_r-p_r\|_1\le E_r,
 \qquad
 \left\|\sum_{h\in(\Z/P\Z)^\times}\widetilde p_{r,h}
          -\frac{|(\Z/P\Z)^\times|}{P}p_r\right\|_1\le E_r.
\]
Consequently the color's total error against the target in
\eqref{eq:signal-refreshed} is bounded by
\begin{equation}
 T_r+(1+2\gamma\|d_r\|_\infty)E_r
       +2\gamma\rho_r(\|p_r\|_1+E_r).
 \label{eq:signal-refresh-error}
\end{equation}
Choose the refresh errors, displacements, and transmission errors
small enough and sum this bound over the finite list.  This proves
\eqref{eq:signal-refreshed}, including the exact value of $\lambda$.
Only residue balancing after summing over a color was needed, because
\eqref{eq:signal-refreshed-match} holds pointwise before the sum.

The remaining assertions follow from the two passage lemmas and
\eqref{eq:signal-index-drift}--\eqref{eq:signal-temporal-drift}.
Specifically, the triangle inequality gives
\[
 \dist_\T(x',z')\le\dist_\T(x,z)+|x'-x|+\dist_\T(z',z).
\]
Positive separation distances and the fixed finite number of passages
allow displacement budgets smaller than all band margins and color
gaps.  Constants such as $\lambda$ are fixed before these error budgets
are imposed.  The required $c$ and passage integers may then be chosen
afresh after each actual finite packet list is known.
\end{proof}
